\chapter{Introducción}

\begin{comment}
Este manual tiene como propósito ser una guía lo más completa posible del material necesario para participar en el Concurso Internacional Universitario de Programación (ICPC, por sus siglas en inglés) en cuánto a la parte matemática se refiere, es cierto que esta competencia es interdisciplinaria y usualmente se busca que los equipos se conformen por gente con un trasfondo fuerte en algoritmia y además en matemáticas.Este libro se enfocará en los conocimientos necesarios de matemáticas que le conviene saber a un concursante por lo que se recomienda para personas estudiantes de matemáticas que quieran iniciarse en la ICPC o personas participantes de ICPC que busquen expandir sus conocimientos matemáticos.

Aunque este manual tenga una gran carga de teoría y definiciones, la familiarización con cada uno de los contenidos de este manual radicará no sólo en leer los capítulos pero también en intentar resolver los problemas y ejercicios sugeridos, si bien algunos problemas fueron recabados de jueces en línea, muchos de ellos fueron diseñados con mucho cariño con la única finalidad de servir para el lector de este manual.

Es importante mencionar que si bien este manual está ordenado por capítulos, estos no están ordenados por dificultad, si no por área de las matemáticas a las que pertenecen, por lo que el lector no necesita necesariamente seguir el orden del indice.
% Esta parte modificarla para explciar un poco más la relevancia de cada area en competencias.
El primer capítulo se presentan técnicas de álgebra avanzada usualmente impartidas dentro de los cursos de Algebra Lineal y Algebra Moderna, el segundo capítulo se enfocará en la parte de teoría de números, se cubrirán conocimientos impartidos en Teoría de Números I y II, así como propiedades básicas usualmente vistas en la Teoría Análitica de Números. La sección de Combinatoria ocupará técnicas y propiedades de combinatoria básicas

La segunda parte de este manual esta dedicada a las soluciones de cada una de los problemas presentados en este libro con explicación y código fuente así como algunos hints para poder avanzar en caso de que el lector se encuentre atorado y no quiera leer la solución. Se recomienda al lector realmente intentar resolver los problemas antes de leer las soluciones, pues la parte más valiosa de este manual no es en sí mismo los temas explicados pero la aplicación de cada uno de ellos a los problemas presentados donde las soluciones puede que no requieran temas difíciles pero sí mucho ingenio.

%Escribir una nota personal de por que quise escribir este manual
\end{comment}

Este manual es una guía de apoyo para fortalecer la preparación matemática en programación competitiva, con enfoque en el Concurso Internacional Universitario de Programación (ICPC). Está dirigido tanto a estudiantes de matemáticas que buscan iniciarse en esta área como a participantes que desean consolidar y ampliar sus herramientas teóricas.
\\\\
Si bien existen libros de programación competitiva ampliamente reconocidos, gran parte de ellos se concentra en algoritmos, implementación y estructuras de datos. Este manual busca complementar ese panorama con un enfoque explícito en la formación matemática que requiere la competencia y con una propuesta adaptada al contexto académico de la Facultad de Ciencias.
\\\\
Además de la exposición conceptual, el manual incluye un conjunto de problemas cuidadosamente seleccionados y también problemas de autoría propia marcados con una estrella ($\filledstar$), diseñados para reforzar de manera directa los conocimientos presentados en cada sección. Para poder acceder y enviar soluciones a los problemas de autoría se creo un grupo del juez Codeforces, al cuál se podrá acceder por medio del \href{https://codeforces.com/contestInvitation/53c4181731f9e66f98b546f60ba5403a60118e4b}{link}\footnote{https://codeforces.com/contestInvitation/53c4181731f9e66f98b546f60ba5403a60118e4b}.
\\\\
Este manual se ha escrito con la idea que los alumnos puedan usarle de referencia, teniendo en cuenta que, en general en la ICPC importa mucho más la practicidad que la formalidad. Sin embargo, este autor es consciente de la importancia que tienen las demostraciones en el mundo de las matemáticas. Tratando de equilibrar estas dos perspectivas, a lo largo del texto se han desarrollado las ideas y conceptos desde una noción intuitiva y se han dado referencias a demostraciones completas para el alumno que decida leerlas.
\\\\
La organización del contenido responde a áreas temáticas, no a niveles de dificultad, con el fin de que cada lector avance de acuerdo con sus intereses y necesidades. En el Capítulo 2 se presenta una introducción general a la competencia ICPC, su dinámica y sus reglas principales; en el Capítulo 3 y 4 se abordan los fundamentos de entrada/salida y el uso de jueces en línea. 
\\\\
Del capítulo 6 al capítulo 9, se desarrollan las áreas de Álgebra, Teoría de números, Combinatoria, Geometría y Teoría de Gráficas. Algunos temas están marcadas por un asterisco (*), indicando que es un tema con mayor dificultad. No todos los temas presentados en este manual están acompañados de un problema motivador propio. Esto no es un descuido, sino una decisión deliberada; algunos temas son herramientas que raramente aparecen de forma aislada, sino como piezas de una solución más grande. Conocer la función $\phi$ de Euler o el lema de Burnside puede no resolver ningún problema por sí solo, pero puede ser exactamente la pieza que faltaba para reducir un problema aparentemente irresoluble a uno manejable.
 
La intención de este manual no es que el lector, al terminarlo, haya agotado el estudio de estas áreas. Las áreas exploradas en este manual son campos vastísimos. Lo que sí se busca es que el lector adquiera un repertorio de herramientas con trasfondo matemático que aparecen con frecuencia en las fechas nacionales y regionales del ICPC, y que entienda cómo se usan, no como recetas, sino como ideas.
 
Quizás lo más valioso que puede ofrecer este manual no es el contenido en sí, sino la forma de pensar que propone: aprender a reconocer patrones, a sospechar cuándo una suma de divisores esconde una convolución o cuándo una trayectoria en un grafo puede traducirse a un sistema de congruencias. Esa capacidad de transformar un problema a una pregunta matemática disfrazada es lo que distingue a un competidor sólido. 
\\\\
Al final del documento se incluyen las editoriales de los problemas presentados. Estas editoriales comienzan con \textit{hints} para orientar el proceso de resolución y no muestran de inmediato la solución completa. Para obtener mejores resultados, se sugiere seguir este orden de trabajo: primero intentar resolver cada problema por cuenta propia, después apoyarse en los \textit{hints} cuando sea necesario y, finalmente, consultar la editorial para comparar enfoques y comprender la solución propuesta.
